\section{考察}
\label{sec:discussion}

\emph{実用上の指針．}
結果から，感情ごとに少数のラベルしかない新しいコーパスに対する4つの推奨が得られる．
第1に，ソースとターゲットの特徴を別々に中心化する．
これはラベルを使わず，ソースのみの分類器がもつ判断基準による損失の大部分を取り除く（\ref{sec:part1-fix}節）．
第2に，すべてのソースから出発し，重みには強いアンカーをかけたまま，緩いアンカーをかけた切片に$K$本のラベルを使う．
ラベルは切片をよく決め，重みをうまく決めないからである（\ref{sec:part2-start}節）．
第3に，どの発話にラベルを付けるかをラベルを使わずに選んでも，よい出発点の向きが直ることは期待しない（\ref{sec:part2-limits}節）．
第4に，ターゲットの発話とともに他のモダリティが記録されているならば，それがラベル付けのときだけであっても，それらのモダリティでソースから学習した先生にプール全体のラベルを付けさせる．
配備する分類器は単一モダリティのままでよい（\ref{sec:part2-crossmodal}節と\ref{sec:comparison}節）．

\emph{テキストのふるまいが異なる理由．}
テキストの特徴は判断基準による損失が小さく，向きによる損失が大きい（表\ref{tab:part1}）．
このため，他の構成で判断基準を直す少数のラベルには，テキストでは直すべき判断基準のずれがほとんど無く，テキストだけの構成では利得のうち判断基準の割合は最大で\NUM{p1-e29-critshare-k1-t}である．
2つのコストによる同じ説明がこの場合にも当てはまる．
利得のうち判断基準から生じる割合は，損失のうち判断基準が引き起こす割合に従う．
提案法では，テキストの生徒も音声の生徒と同じくらい融合の先生から利得を得る．
融合の先生が，テキストのラベルに欠けている向きの情報を補うからである．

\emph{エントロピー最小化とクラスタの仮定．}
エントロピー最小化は判定の境界をラベルのないデータの密な領域から遠ざけ，クラスが分かれたクラスタをなすときに役立つ．
本論文のデータでは，エントロピー最小化は，AVTとATの構成のIEMOCAPの即興のセッションでは役立ち，MELDとCMU-MOSEI，およびVTとAVの構成の同じIEMOCAPのターゲットでは大きく失敗する（\ref{sec:comparison-fixed}節）．
別のモダリティからの弱いラベルは，モダリティ間の相補性という別の性質に依拠しており，その利得のばらつきはずっと小さい．
$K=0$と$K=5$のターゲットとセルの\NUM{st-b-tgt-n}の組み合わせのうち，提案法の拡張版が出発点を下回るのは\NUM{st-b-tgtneg-n}で，その差は最大で\NUM{st-b-tgtneg-worst}ポイントである．

\emph{ソースとは何か．}
第I部はソースの価値（value of a source）を向きのラベルの本数で与え，第II部は先生の価値が出発点に対する先生のUARの優位に応じて増減することを示す．
どちらの量も配備の前に推定できる．
$q$と$e$はラベル付きのソースと仮のターゲットから，$J$は仮のターゲットから推定する．
これらの量は，ターゲットのラベルを増やすこと，ソースを加えること，ラベル付けのときにモダリティを加えることのどれがよりよい投資かを示す．
